\subsection{DDA com \textit{Monte Carlo Tree Search} e o estado do jogador}    
    O método que utiliza um agente de DDA baseado em 
    \textit{Monte Carlo tree search} (MCTS) em um jogo
    de luta é apresentado em 
    \cite{demediuk2017monte,ishihara2018monte}. Nesse 
    método, a dificuldade do jogo é manipulada por meio
    da implementação de um agente de DDA que utiliza a 
    diferença de \textit{health points} (HP) entre o 
    jogador e o agente como função de pontuação, de 
    maneira semelhante à abordagem apresentada em 
    \cite{moon2022diversifying}.
    
    Contudo, essas técnicas de medição das habilidades dos
    jogadores, que utilizam funções de pontuação heurísticas
    especificadas pelos desenvolvedores, pressupõem que o
    ajuste da dificuldade do jogo por meio dessas métricas
    (com base em HP, pontuação no jogo, dano, etc.)
    pode aumentar o \textit{player experience} (PX). 
    No entanto, a PX e a dificuldade do jogo 
    ajustada por essas medidas frequentemente apresentam 
    inconsistências
    \cite{ang2017comparing,constant2019dynamic}.

    Por isso, em \cite{moon2022diversifying}, os estados 
    do jogador são utilizados no treinamento da IA do 
    jogo juntamente com o MCTS, permitindo que o DDA se 
    ajuste automaticamente e progressivamente às ações 
    dos oponentes. Essa abordagem considera que a 
    dificuldade do jogo pode evocar diversos estados 
    cognitivos e emocionais que podem impactar a PX 
    \cite{kivikangas2011review}. Nesse contexto, o MCTS 
    busca a próxima ação, considerada quase 
    ótima para a IA do jogo, por meio da simulação de 
    um grande número de situações possíveis a 
    partir do estado atual do jogo. 
    
    Para tanto, as simulações do MCTS consistem em quatro 
    etapas: \textit{selection}, \textit{expansion}, 
    \textit{simulation} e \textit{backpropagation}. 
    Em cada iteração, o estado atual é definido como o nó
    raiz. Em seguida, as quatro etapas da busca em 
    árvore são executadas dentro de um limite
    de tempo predeterminado. Após a conclusão da busca, 
    o agente escolhe o nó filho mais visitado ou aquele 
    que apresenta o maior valor para realizar a próxima 
    ação \cite{moon2022diversifying}.

    Sendo assim, a principal diferença entre o modelo 
    proposto em \cite{moon2022diversifying} e o método de 
    DDA heurístico é a utilização direta do estado do 
    jogador como função de pontuação para o MCTS, em vez
    de funções heurísticas, como pontuação do jogo, taxa
    de vitória e diferença de HP. O modelo do estado do
    jogador combina registros de partidas reais e de 
    partidas simuladas pelo MCTS para prever como o 
    estado do jogador se modificará no futuro por meio
    de técnicas de \textit{machine learning} (ML) 
    \cite{moon2022diversifying}.